\documentclass[11pt,a4paper]{article}

\usepackage[margin=25mm]{geometry}
\usepackage[T1]{fontenc}
\usepackage[utf8]{inputenc}
\usepackage{newtxtext,newtxmath}
\usepackage{amsmath,bm}
\usepackage{graphicx}
\usepackage{booktabs}
\usepackage{xcolor}
\usepackage{microtype}
\usepackage{xurl}
\usepackage[hidelinks]{hyperref}

\setcounter{secnumdepth}{3}
\setlength{\parindent}{2em}
\setlength{\parskip}{0.25em}
\setlength{\emergencystretch}{2em}
\allowdisplaybreaks
\numberwithin{equation}{section}

% Revision markup. Delete these two macros and their uses before submission.
\newcommand{\rev}[1]{\textcolor{red}{#1}}
\newcommand{\todo}[1]{\textcolor{red}{\textbf{[TODO: #1]}}}


\title{On Craig--Bampton Model Reduction for Multi-Degree-of-Freedom Coupled Real-Time Hybrid Simulation\\[4pt]
\large Draft with Sections 2 to 5}
\author{Ning Li \and Yu Liang}
\date{}

\begin{document}

\maketitle

\begin{abstract}
	
\end{abstract}

\section{Introduction} \label{sec1}

Experimental assessment of structural systems under earthquake loading requires the reproduction of inertia interaction, restoring-force nonlinearity, and loading-rate effects. Shake-table testing reproduces the complete dynamic response of a specimen, but the payload, stroke, and laboratory space restrict the size and configuration of the test structure. Pseudodynamic testing reduces this demand by computing the inertia and damping forces and measuring the restoring force of a specimen loaded step by step \cite{Mahin1989}\,[DOI \nolinkurl{10.1061/(ASCE)0733-9445(1989)115:8(2113)}]. Its loading rate remains below the seismic rate and does not reproduce the response of rate-dependent components. Real-time hybrid simulation (RTHS) advances the numerical solution and the physical loading on the same time scale. The structure is partitioned into a numerical substructure and a physical substructure, and displacement and restoring force are exchanged within each integration step \cite{Nakashima1992}\,[DOI \nolinkurl{10.1002/eqe.4290210106}]. This combination retains the experimental representation of nonlinear or rate-dependent components without requiring a full-structure dynamic test \cite{Gao2013}\,[DOI \nolinkurl{10.1002/eqe.2246}].

The extension from a single imposed degree of freedom (DOF) to multi-DOF and multi-axial RTHS has made the transfer system a coupled multiple-input multiple-output plant. Model-based multi-actuator control was developed to coordinate several loading channels \cite{Phillips2013}\,[DOI \nolinkurl{10.1061/(ASCE)EM.1943-7889.0000493}]. Multi-axial frameworks and load and boundary condition boxes subsequently addressed actuator coupling, nonlinear kinematics, and control-structure interaction \cite{Fermandois2017}\,[DOI \nolinkurl{10.1007/s11803-017-0407-8}] \cite{Najafi2020}\,[DOI \nolinkurl{10.1016/j.engstruct.2020.110868}]. A state-of-the-art review and a common benchmark consolidated these developments and enabled direct comparison of advanced compensation and control methods \cite{Najafi2023}\,[DOI \nolinkurl{10.1016/j.engstruct.2022.115284}] \cite{Condori2023}\,[DOI \nolinkurl{10.3389/fbuil.2023.1270996}]. Adaptive, decentralized, and data-driven strategies have further improved the tracking of the prescribed channels \cite{Quiroz2024}\,[DOI \nolinkurl{10.3389/fbuil.2024.1394952}] \cite{RuizSong2024}\,[DOI \nolinkurl{10.3389/fbuil.2024.1477804}] \cite{Xu2024}\,[DOI \nolinkurl{10.3389/fbuil.2024.1374819}]. These methods begin after the actuated DOFs have been selected. The number of independent kinematic conditions that can be imposed remains bounded by the available actuators, reaction system, and hydraulic capacity.

An MDOF physical substructure can therefore contain more DOFs that should be imposed than the transfer system can actuate independently. Increasing the number of actuators raises the hardware demand and enlarges the coupled control problem. Condensation provides a different solution by representing non-actuated DOFs through a smaller set of retained DOFs that the transfer system imposes. In this setting, model reduction does more than lower the order of the equations. It determines how every condensed DOF is reconstructed from the actuator-driven DOFs and therefore changes the path by which actuator errors enter the RTHS loop.

Guyan static condensation and Craig--Bampton component-mode synthesis provide two established bases for this operation. Guyan condensation recovers the condensed DOFs from a static relation and produces a compact model \cite{Guyan1965}\,[DOI \nolinkurl{10.2514/3.2874}]. Craig--Bampton condensation augments the same constraint contribution with selected fixed-interface modes \cite{CraigBampton1968}\,[DOI \nolinkurl{10.2514/3.4741}]. Model reduction has already been introduced into hybrid simulation to reduce the frequency bandwidth of stiff hybrid models \cite{Miraglia2020}\,[DOI \nolinkurl{10.1002/eqe.3262}]. Data-driven surrogates, dynamic condensation, and proper orthogonal decomposition have also been developed to satisfy the real-time computation limit of high-dimensional numerical substructures \cite{Tsokanas2022}\,[DOI \nolinkurl{10.1016/j.mechmachtheory.2022.105072}] \cite{Mucha2023}\,[DOI \nolinkurl{10.1007/s11012-023-01675-0}] \cite{Zhang2024}\,[DOI \nolinkurl{10.1002/eqe.4221}]. These studies establish the value of reduction for bandwidth control and computational feasibility. They do not connect an actuation-limited reduction of the physical substructure to channel-wise delay propagation and closed-loop delay margin.

Actuator delay forms a second well-developed line of RTHS research. The phase lag of the measured restoring force can supply energy to the loop and act as negative damping \cite{Horiuchi1999}\,[DOI \nolinkurl{10.1002/(SICI)1096-9845(199910)28:10\textless1121::AID-EQE858\textgreater3.0.CO;2-O}]. Delay-differential-equation models established continuous-time stability boundaries \cite{Wallace2005}\,[DOI \nolinkurl{10.1002/eqe.513}]. Stability analysis was then extended to explicit integration schemes \cite{ChenRicles2008}\,[DOI \nolinkurl{10.1002/eqe.775}], multiple delay sources \cite{Mercan2008}\,[DOI \nolinkurl{10.1002/eqe.814}], and MDOF systems through discrete-time root-locus analysis \cite{Zhu2015}\,[DOI \nolinkurl{10.1002/eqe.2467}]. Recent formulations have treated coupled nonlinear behavior and delay-induced spectral and energy redistribution \cite{Huang2022}\,[DOI \nolinkurl{10.1002/eqe.3667}] \cite{Huang2025}\,[DOI \nolinkurl{10.1002/eqe.70046}]. These formulations specify the delayed DOFs and their coefficient matrices before the stability analysis. They do not identify how the reduction basis selects the part of the physical-substructure response exposed to each actuator delay.

Condensing the loaded side of the interface creates this missing connection. A condensed DOF is neither commanded nor measured, and its response is recovered from the retained DOFs through the transformation matrix. The retained physical DOFs are actuator driven and therefore delayed. Under Guyan condensation, the reconstructed response consists only of combinations of these delayed quantities. Under Craig--Bampton condensation, the fixed-interface modal coordinates are solved inside the reduced model and carry no explicit actuator-delay operator. The reduction basis consequently determines the fraction of the physical-substructure response that enters the closed loop through delayed channels. Accuracy and stability must therefore be assessed together.

This paper establishes a delay-dependent reduction framework for an MDOF RTHS whose physical substructure contains more imposed DOFs than the transfer system can actuate independently. The Guyan and Craig--Bampton transformations are formulated on a common retained and condensed partition, and the actuator delays are propagated through each transformation channel by channel. The reduced numerical and physical substructures are assembled into a discrete-time closed loop with independent integer channel delays and a linear quadratic regulator designed from the corresponding reduced model. The admissible delay pairs are obtained from the modulus of the dominant closed-loop pole. A modal redistribution ratio is further defined from the full-order and reduced modal solutions to quantify the low-order modal content moved onto the actuator-driven DOFs before a delay scan is performed. Two divisions of a three-storey three-bay frame, including one in which a principal storey displacement is condensed, establish the accuracy, stability, and range of application of the two methods.

\section{Structural modeling and condensation formulations} \label{sec2}

In a multi-degree-of-freedom (MDOF) real-time hybrid simulation (RTHS), the \rev{physical substructure often contains} more interface coordinates than \rev{the transfer system can impose independently. Model reduction is adopted to express the coordinates without independent actuation through} a reduced set of experimentally imposed coordinates.
In this section, the structural modeling framework and the condensation formulations are established. The full-order equation of motion is \rev{decomposed} into the contributions of the numerical and physical substructures. Guyan static condensation and Craig--Bampton dynamic condensation are then formulated on a common partition of retained and condensed coordinates, and the two methods differ only in the \rev{construction of the transformation matrix}.

\subsection{Equations of motion and substructure partitioning}
\rev{The equation of motion of the full-order structure under a ground acceleration is expressed as}
\begin{equation}
	\mathbf{M}\ddot{\bm{x}}(t)+\mathbf{C}\dot{\bm{x}}(t)+\mathbf{K}\bm{x}(t)
	=-\mathbf{M}\bm{\Gamma}a_{\mathrm{g}}(t),
	\label{eq:full-order-eom}
\end{equation}
where \(\mathbf{M}\), \(\mathbf{C}\), and \(\mathbf{K}\) are the mass, damping, and stiffness matrices of the full-order structure\rev{, respectively, \(\bm{x}(t)\) is the generalized displacement vector of the full-order structural DOFs measured relative to the ground}, \(\bm{\Gamma}\) is the influence vector, and \(a_{\mathrm{g}}(t)\) is the ground acceleration.

After the structure is divided into numerical and physical substructures, the corresponding equations are written as
\begin{subequations}
	\label{eq:substructure-eom}
	\begin{align}
		\mathbf{M}_{\mathrm{n}}\ddot{\bm{x}}_{\mathrm{n}}
		+\mathbf{C}_{\mathrm{n}}\dot{\bm{x}}_{\mathrm{n}}
		+\mathbf{K}_{\mathrm{n}}\bm{x}_{\mathrm{n}}
		=\bm{f}_{\mathrm{n}}{}+\bm{\lambda}_{\mathrm{n}},
		\label{eq:numerical-eom}\\
		\mathbf{M}_{\mathrm{e}}\ddot{\bm{x}}_{\mathrm{e}}
		+\mathbf{C}_{\mathrm{e}}\dot{\bm{x}}_{\mathrm{e}}
		+\mathbf{K}_{\mathrm{e}}\bm{x}_{\mathrm{e}}
		=\bm{f}_{\mathrm{e}}{}+\bm{\lambda}_{\mathrm{e}},
		\label{eq:physical-eom}
	\end{align}
\end{subequations}
where \(\bm{x}_{\mathrm{n}}\) and \(\bm{x}_{\mathrm{e}}\) \rev{are} the DOFs of the numerical and physical substructures, respectively\rev{, \(\bm{f}_{\mathrm{n}}\) and \(\bm{f}_{\mathrm{e}}\) are the external load vectors carrying the share of the seismic inertia load of Eq. \eqref{eq:full-order-eom} taken by each substructure, and \(\bm{\lambda}_{\mathrm{n}}\) and \(\bm{\lambda}_{\mathrm{e}}\) are the interface force vectors acting on the two substructures. The interface coordinates are shared by the two substructures and the interface forces balance}, \(\bm{\lambda}_{\mathrm{n}}+\bm{\lambda}_{\mathrm{e}}=\bm{0}\). \rev{The partition assigns one part of the structure to numerical solution and the other to physical implementation, and it does not dynamically decouple the two substructures.}

For a condensation operation, the DOFs of a substructure are arranged into a retained set \(\bm{u}_{\mathrm{r}}\) and a condensed set \(\bm{u}_{\mathrm{c}}\), and \rev{the} equation of motion takes the form
\begin{equation}
	\begin{bmatrix}
		\mathbf{M}_{\mathrm{rr}} & \mathbf{M}_{\mathrm{rc}}\\
		\mathbf{M}_{\mathrm{cr}} & \mathbf{M}_{\mathrm{cc}}
	\end{bmatrix}
	\begin{bmatrix}
		\ddot{\bm{u}}_{\mathrm{r}}\\
		\ddot{\bm{u}}_{\mathrm{c}}
	\end{bmatrix}
	+
	\begin{bmatrix}
		\mathbf{C}_{\mathrm{rr}} & \mathbf{C}_{\mathrm{rc}}\\
		\mathbf{C}_{\mathrm{cr}} & \mathbf{C}_{\mathrm{cc}}
	\end{bmatrix}
	\begin{bmatrix}
		\dot{\bm{u}}_{\mathrm{r}}\\
		\dot{\bm{u}}_{\mathrm{c}}
	\end{bmatrix}
	+
	\begin{bmatrix}
		\mathbf{K}_{\mathrm{rr}} & \mathbf{K}_{\mathrm{rc}}\\
		\mathbf{K}_{\mathrm{cr}} & \mathbf{K}_{\mathrm{cc}}
	\end{bmatrix}
	\begin{bmatrix}
		\bm{u}_{\mathrm{r}}\\
		\bm{u}_{\mathrm{c}}
	\end{bmatrix}
	=
	\begin{bmatrix}
		\bm{f}_{\mathrm{r}}\\
		\bm{f}_{\mathrm{c}}
	\end{bmatrix},
	\label{eq:partitioned-local-eom}
\end{equation}
\rev{where} \(\bm{f}_{\mathrm{r}}\) and \(\bm{f}_{\mathrm{c}}\) are the generalized forces acting on the retained and \rev{condensed} coordinates, which \rev{collect the seismic load carried by the substructure together with the interface force}.

The condensed coordinates are expressed through the retained coordinates by a transformation matrix, and the substructure matrices are projected onto the reduced subspace by a congruent transformation,
\begin{equation}
	\begin{gathered}
		\bm{u}(t)=\mathbf{T}\bm{q}(t),\\
		\mathbf{M}_{\mathrm{red}}=\mathbf{T}^{\mathsf{T}}\mathbf{M}\mathbf{T},\quad
		\mathbf{C}_{\mathrm{red}}=\mathbf{T}^{\mathsf{T}}\mathbf{C}\mathbf{T},\quad
		\mathbf{K}_{\mathrm{red}}=\mathbf{T}^{\mathsf{T}}\mathbf{K}\mathbf{T},\quad
		\bm{f}_{\mathrm{red}}=\mathbf{T}^{\mathsf{T}}\bm{f},
	\end{gathered}
	\label{eq:general-projection}
\end{equation}
where \(\bm{q}(t)\) is the reduced generalized-coordinate vector \rev{and \(\mathbf{T}\) is the transformation matrix}. Guyan and Craig--Bampton condensation share Eq. \eqref{eq:general-projection} and differ only in the \rev{construction of \(\mathbf{T}\)}.

The two substructures are condensed separately and are \rev{coupled at the interface coordinates retained by both reduced models}, where the interface forces cancel. The assembled reduced model is governed by
\begin{equation}
	\mathbf{M}_{\mathrm{red}}\ddot{\bm{q}}(t)
	+\mathbf{C}_{\mathrm{red}}\dot{\bm{q}}(t)
	+\mathbf{K}_{\mathrm{red}}\bm{q}(t)
	=\bm{f}_{\mathrm{red}}(t),
	\label{eq:assembled-reduced}
\end{equation}
\rev{where} \(\bm{q}(t)\) collects the generalized coordinates of the two reduced substructures, with the shared interface coordinates counted once. Only the interface coordinates \rev{retained by both reduced models} are matched, \rev{since} these are the coordinates that the transfer system can impose. \rev{A geometric interface coordinate outside the coupling set carries neither the compatibility condition nor the equilibrium condition of Eq. \eqref{eq:substructure-eom}, and each reduced substructure recovers it from its own transformation. The same restriction applies to the test itself. Once the interface carries more coordinates than there are actuation channels}, the specimen responds on the unactuated interface coordinates according to its own dynamics, and no loading system can enforce the numerical model there.

\subsection{Guyan static condensation} \label{sec21}

Guyan condensation derives the condensed coordinates from static equilibrium \cite{Guyan1965}\,\rev{[DOI \nolinkurl{10.2514/3.2874}]}. The reduction basis is built from the static response of the condensed coordinates\rev{. Their inertia, their damping and the load acting on them are therefore dropped} from the lower block of Eq. \eqref{eq:partitioned-local-eom}, which leaves \(\mathbf{K}_{\mathrm{cr}}\bm{u}_{\mathrm{r}}+\mathbf{K}_{\mathrm{cc}}\bm{u}_{\mathrm{c}}=\bm{0}\). The condensed coordinates are \rev{approximated} by
\begin{equation}
	\bm{u}_{\mathrm{c}}=\bm{\Psi}_{\mathrm{G}}\bm{u}_{\mathrm{r}},
	\qquad
	\bm{\Psi}_{\mathrm{G}}=-\mathbf{K}_{\mathrm{cc}}^{-1}\mathbf{K}_{\mathrm{cr}},
	\label{eq:guyan-coordinate-relation}
\end{equation}
\rev{where \(\bm{\Psi}_{\mathrm{G}}\) is the static condensation matrix, which requires \(\mathbf{K}_{\mathrm{cc}}\) to be non-singular.}
The Guyan transformation is
\begin{equation}
	\bm{u}=\mathbf{T}_{\mathrm{G}}\bm{q}_{\mathrm{G}},
	\qquad
	\mathbf{T}_{\mathrm{G}}=
	\begin{bmatrix}
		\mathbf{I}\\
		\bm{\Psi}_{\mathrm{G}}
	\end{bmatrix},
	\qquad
	\bm{q}_{\mathrm{G}}=\bm{u}_{\mathrm{r}},
	\label{eq:guyan-transformation}
\end{equation}
\rev{where \(\bm{q}_{\mathrm{G}}\) is the reduced coordinate vector, which contains the retained coordinates alone.} The reduced matrices follow from Eq. \eqref{eq:general-projection} with \(\mathbf{T}=\mathbf{T}_{\mathrm{G}}\), and the \rev{response} of the condensed coordinates is recovered \rev{from} Eq. \eqref{eq:guyan-coordinate-relation}.

The loads \rev{on} the condensed coordinates, \rev{including} their share of the seismic inertia load and the interface force, are \rev{projected} through \(\mathbf{T}_{\mathrm{G}}^{\mathsf{T}}\bm{f}\)\rev{, and they no longer act on independent coordinates. Guyan condensation is accordingly most accurate when the condensed coordinates follow the retained coordinates through quasi-static deformation, and it loses accuracy once the inertia of the condensed coordinates contributes to the response.}

\subsection{Craig--Bampton dynamic condensation}
\label{sec:cb-condensation}

Craig--Bampton condensation keeps the retained coordinates explicit and represents the condensed coordinates through constraint modes and selected fixed-interface modes \cite{CraigBampton1968}\,\rev{[DOI \nolinkurl{10.2514/3.4741}]}. \rev{The experimentally imposed coordinates form the retained set of the physical substructure, and internal master coordinates are additionally retained in the numerical substructure. An interface coordinate driven by no actuator is therefore condensed, whereas the classical partition keeps every interface coordinate explicit} \cite{Krattiger2019}\,\rev{[DOI \nolinkurl{10.1016/j.ymssp.2018.05.031}]}.

\rev{The dynamic modes are obtained with the retained coordinates constrained, and satisfy}
\begin{equation}
	\mathbf{K}_{\mathrm{cc}}\bm{\Phi}
	=\mathbf{M}_{\mathrm{cc}}\bm{\Phi}\bm{\Omega}^{2},
	\qquad
	\bm{\Phi}^{\mathsf{T}}\mathbf{M}_{\mathrm{cc}}\bm{\Phi}=\mathbf{I},
	\label{eq:fixed-interface-modes}
\end{equation}
where the columns of \(\bm{\Phi}\) are the \rev{fixed-interface modes of the partition and} \(\bm{\Omega}^{2}\) is the diagonal matrix of the corresponding eigenvalues.

The constraint modes describe the static deformation of the condensed coordinates produced by unit displacements of the retained coordinates, which is the relation \(\bm{\Psi}_{\mathrm{G}}\) of Eq. \eqref{eq:guyan-coordinate-relation}. The Craig--Bampton transformation is therefore
\begin{equation}
	\begin{bmatrix}
		\bm{u}_{\mathrm{r}}\\
		\bm{u}_{\mathrm{c}}
	\end{bmatrix}
	=
	\begin{bmatrix}
		\mathbf{I} & \mathbf{0}\\
		\bm{\Psi}_{\mathrm{G}} & \bm{\Phi}_{d}
	\end{bmatrix}
	\begin{bmatrix}
		\bm{u}_{\mathrm{r}}\\
		\bm{\eta}
	\end{bmatrix}
	=
	\mathbf{T}_{\mathrm{CB}}\bm{q}_{\mathrm{CB}},
	\label{eq:cb-transformation}
\end{equation}
where \rev{\(\bm{\Phi}_{d}\) collects the \(d\) fixed-interface modes of lowest frequency, which are the modes retained in the reduced model, and} \(\bm{\eta}\) contains the \rev{corresponding modal coordinates}. The reduced matrices follow from Eq. \eqref{eq:general-projection} with \(\mathbf{T}=\mathbf{T}_{\mathrm{CB}}\).

The response of the condensed coordinates is approximated as
\begin{equation}
	\bm{u}_{\mathrm{c}}(t)=
	\bm{\Psi}_{\mathrm{G}}\bm{u}_{\mathrm{r}}(t)
	+\bm{\Phi}_{d}\bm{\eta}(t).
	\label{eq:cb-internal-recovery}
\end{equation}
The first term is the static contribution of Guyan condensation\rev{, and the second term retains selected internal vibration patterns, which is the only difference between the two transformations}.


\section{Delay-dependent stability formulation} \label{sec3}

In an MDOF RTHS, each actuated coordinate of the physical substructure is driven by a separate actuator, and each actuator introduces its own response delay. \rev{Condensation expresses the condensed coordinates through the retained coordinates, and the delays acting on the retained coordinates are therefore carried by} the reconstructed response of the condensed coordinates. In this section, the delay-dependent stability formulation of the condensed RTHS system is established. The actuator delays are \rev{represented} as integer multiples of the integration time step, and \rev{their transmission through} the Guyan and Craig--Bampton transformations is identified. The condensed substructures are then assembled into a discrete-time closed-loop model with an LQR controller, and \rev{the} stability is assessed from the modulus of the dominant pole \rev{of the closed loop}.

\subsection{Multiple-delay representation and delay propagation through condensation}
\label{sec:multiple-delay}

The delay of an RTHS system originates from the numerical integration, the data exchange between the numerical and physical substructures, the conversion between digital and analogue signals, and the response of the servo-hydraulic actuator \cite{Carrion2007}. \rev{The actuator response delay is retained in the present formulation and the other three sources are excluded, which isolates the interaction between the delay and the model reduction.}

The actuator delay is taken as an integer multiple of the integration time step \cite{Zhu2015}\,\rev{[DOI \nolinkurl{10.1002/eqe.2467}]}. \rev{The delay then becomes a power of \(z\) after the \(z\) transform, which avoids the transcendental term produced by a continuous-time delay and admits several delay channels at the same time.} The delay of the \(\ell\)th actuator and the displacement that this actuator imposes on the specimen are
\begin{equation}
	\tau_{\ell}=n_{\ell}\Delta t,
	\qquad
	\hat{u}_{\mathrm{r},\ell}(t)=u_{\mathrm{r},\ell}(t-\tau_{\ell}),
	\qquad
	\ell=1,\dots,n_{\mathrm{a}},
	\label{eq:integer-delay}
\end{equation}
where \(\Delta t\) is the integration time step, \(n_{\ell}\) is a non-negative integer\rev{, with \(n_{\ell}=0\) denoting a delay-free channel}, \(n_{\mathrm{a}}\) is the number of actuators, \(u_{\mathrm{r},\ell}\) is the retained coordinate driven by the \(\ell\)th actuator, and \(\hat{u}_{\mathrm{r},\ell}\) is the displacement \rev{imposed} on the specimen. Each actuator is assigned its own value of \(n_{\ell}\), \rev{which represents channels of different equivalent pure delay}. A pure delay reproduces the phase lag of a channel \rev{rather than} its bandwidth or its amplitude attenuation. The closed loop \rev{carrying} these delays is shown in Fig. \ref{fig:rths-loop}. The numerical substructure supplies the target displacement of each actuated coordinate, the transfer system imposes it on the physical substructure with the delay of Eq. \eqref{eq:integer-delay}, and the measured displacement and \rev{restoring force are returned to the numerical substructure and to the regulator acting on the actuated coordinates}.

\begin{figure}[htbp]
	\centering
	% Selected option: Figure 1-D. Compare with TU thesis Fig. 2-1 (PDF p.21).
	\includegraphics[width=0.96\textwidth]{figure/selected_v2/fig01_rths_loop_optionD.png}
	\caption{Closed loop of an MDOF RTHS with multiple actuators.}
	\label{fig:rths-loop}
\end{figure}

The retained coordinates of the physical substructure are the only coordinates that the transfer system imposes, and each of them is driven by one actuator. The condensed coordinates are neither commanded nor measured, and their response is reconstructed from the recovery relation of the condensation. The two condensation methods reconstruct this response in different ways.

The reconstructed response follows from Eqs. \eqref{eq:guyan-coordinate-relation} and \eqref{eq:cb-internal-recovery} as
\begin{subequations}
	\label{eq:delay-recovery}
	\begin{align}
		\bm{u}_{\mathrm{c}}^{\mathrm{rec}}(t)&=\bm{\Psi}_{\mathrm{G}}\hat{\bm{u}}_{\mathrm{r}}(t),
		\label{eq:guyan-delay-recovery}\\
		\bm{u}_{\mathrm{c}}^{\mathrm{rec}}(t)&=\bm{\Psi}_{\mathrm{G}}\hat{\bm{u}}_{\mathrm{r}}(t)
		+\bm{\Phi}_{d}\bm{\eta}(t),
		\label{eq:cb-delay-recovery}
	\end{align}
\end{subequations}
for Guyan and \rev{Craig--Bampton condensation, respectively}. Every entry of \(\hat{\bm{u}}_{\mathrm{r}}\) is a delayed quantity\rev{. The Guyan reconstruction is therefore built from delayed quantities alone, and} \(\bm{\Psi}_{\mathrm{G}}\) redistributes the delayed channels over all condensed coordinates. The fixed-interface modal coordinates \(\bm{\eta}\) are solved within the reduced model and do not pass through the transfer system\rev{. The second term of Eq. \eqref{eq:cb-delay-recovery} accordingly carries no explicit delay operator, and it reaches the delayed motion of the retained coordinates only through the coupling of the reduced equations}. For the same actuated coordinates and the same actuator delays, the delayed channels \rev{reach} the reconstructed response through the constraint-mode term alone under Craig--Bampton condensation, and through every term under Guyan condensation.

The formulation is developed for a virtual RTHS, \rev{in which} \(\bm{\eta}\) of the physical substructure is a state of the simulated specimen model.

\subsection{Discrete-time closed-loop pole criterion with LQR control}
\label{sec:pole-criterion}

\rev{The assembled reduced model of Eq. \eqref{eq:assembled-reduced} is taken as the starting point. The actuator delay affects only the terms acting on the coordinates imposed by the transfer system, and each channel carries its own delay. The assembled damping and stiffness matrices are therefore split channel by channel, and the equation of motion of the assembled reduced system is}
\begin{equation}
	\mathbf{M}_{\mathrm{red}}\ddot{\bm{q}}(t)+\mathbf{C}_{0}\dot{\bm{q}}(t)+\mathbf{K}_{0}\bm{q}(t)
	+\sum_{\ell=1}^{n_{\mathrm{a}}}
	\left[
	\mathbf{C}_{\ell}\dot{\bm{q}}(t-\tau_{\ell})
	+\mathbf{K}_{\ell}\bm{q}(t-\tau_{\ell})
	\right]
	=\bm{f}_{\mathrm{red}}(t),
	\label{eq:delayed-eom}
\end{equation}
where \(\bm{q}(t)\) collects the \(n_{q}\) generalized coordinates of the assembled reduced model, \(\mathbf{C}_{\ell}\) and \(\mathbf{K}_{\ell}\) are of the same size as \(\mathbf{C}_{\mathrm{red}}\) and \(\mathbf{K}_{\mathrm{red}}\) and retain only the columns of the reduced physical substructure that act on the coordinate driven by the \(\ell\)th actuator, the remaining columns being zero, and \(\mathbf{C}_{0}=\mathbf{C}_{\mathrm{red}}-\sum_{\ell}\mathbf{C}_{\ell}\) and \(\mathbf{K}_{0}=\mathbf{K}_{\mathrm{red}}-\sum_{\ell}\mathbf{K}_{\ell}\) contain the remaining terms. This split makes the difference \rev{described in} Section \ref{sec:multiple-delay} explicit. \rev{The reduced physical substructure of the Guyan model retains the actuated coordinates alone, and every one of its columns is therefore delayed. The Craig--Bampton model also carries fixed-interface modal coordinates driven by no actuator, and the columns associated with them remain in \(\mathbf{C}_{0}\) and \(\mathbf{K}_{0}\) without delay.}

\rev{The force returned to the loop by the physical substructure is its elastic and damping reaction to the motion imposed by the transfer system. Each actuated channel therefore enters Eq. \eqref{eq:delayed-eom} with its own delay as a scalar factor, and the delays remain independent of each other.} The inertia of the physical substructure is advanced numerically within the assembled state, \rev{and} the assembled mass matrix carries no delay factor. \rev{A term \(\mathbf{M}_{\ell}\ddot{\bm{q}}(t-\tau_{\ell})\) would be added to Eq. \eqref{eq:delayed-eom} once the returned force contained the inertia produced by the delayed motion.}

The assembled system of Eq. \eqref{eq:delayed-eom} is advanced with the CR integration algorithm \cite{ChenRicles2008b}\,\rev{[DOI \nolinkurl{10.1061/(ASCE)0733-9399(2008)134:8(676)}]}, and the sampling period \rev{equals} the integration time step. \rev{A quantity written with the argument \(z\) denotes the \(z\) transform of the corresponding time-domain quantity in the following. The displacement and velocity recursions of the algorithm share one matrix coefficient. Eliminating the acceleration between them and taking the \(z\) transform relates the velocity and the acceleration to the displacement by}
\begin{equation}
	\begin{gathered}
		\dot{\bm{q}}(z)=\frac{z-1}{z\Delta t}\bm{q}(z),\\
		\mathbf{M}_{\mathrm{red}}\ddot{\bm{q}}(z)=
		\left(\mathbf{M}_{\mathrm{red}}
		+\tfrac{1}{2}\Delta t\,\mathbf{C}_{\mathrm{red}}
		+\tfrac{1}{4}\Delta t^{2}\mathbf{K}_{\mathrm{red}}\right)
		\frac{(z-1)^{2}}{z\Delta t^{2}}\bm{q}(z),
	\end{gathered}
	\label{eq:z-derivatives}
\end{equation}
\rev{where} the bracket is the mass term of the CR algorithm\rev{, formed} from the assembled matrices before the channel split of Eq. \eqref{eq:delayed-eom}, since that split describes the signal path of the interface restoring force \rev{rather than} the time integration. A delay of \(n_{\ell}\) steps becomes the scalar factor \(z^{-n_{\ell}}\). The delayed terms are taken from the stored response history\rev{, and the recursion therefore remains explicit in them}.

A linear quadratic regulator acts on the actuated coordinates and supplies a generalized force that reaches the specimen through the actuator channels. Its design model is the assembled reduced model statically condensed onto the actuated coordinates through Eqs. \eqref{eq:guyan-coordinate-relation} and \eqref{eq:guyan-transformation}, with the regulator force as its input. Minimizing a quadratic cost of the design-model states and of the control input, weighted by \(\mathbf{Q}\) and \(\mathbf{R}\), gives the control law
\begin{equation}
	\bm{f}_{\mathrm{L}}(t)=
	-\mathbf{K}_{x}\mathbf{S}_{\mathrm{a}}\bm{q}(t)
	-\mathbf{K}_{v}\mathbf{S}_{\mathrm{a}}\dot{\bm{q}}(t),
	\label{eq:lqr-law}
\end{equation}
\rev{where} \(\mathbf{S}_{\mathrm{a}}\) is the Boolean matrix of size \(n_{\mathrm{a}}\times n_{q}\) that extracts the actuated coordinates from \(\bm{q}\), and \(\mathbf{K}_{x}\) and \(\mathbf{K}_{v}\) are the \(n_{\mathrm{a}}\times n_{\mathrm{a}}\) displacement and velocity parts of the gain.

The regulator output is a generalized force\rev{. Equation \eqref{eq:lqr-law} therefore acts} on the assembled equation as an equivalent feedback stiffness \(\mathbf{K}_{x}\) and an equivalent feedback damping \(\mathbf{K}_{v}\) on the actuated coordinates\rev{, and this action reaches the specimen through the actuator channels of Fig. \ref{fig:rths-loop} with the same channel delays as the imposed displacement. The design model has \(2n_{\mathrm{a}}\) states against \(2n_{q}\) of the assembled model, and Eq. \eqref{eq:lqr-law} is accordingly an output feedback on the actuated coordinates rather than a full-state regulator of the assembled model. Each model supplies its own design model, and equal \(\mathbf{Q}\) and \(\mathbf{R}\) therefore give different gains. A stability domain computed in this way carries the regulator produced by the model itself.}
Substituting Eqs. \eqref{eq:z-derivatives} and \eqref{eq:lqr-law} into Eq. \eqref{eq:delayed-eom} gives the closed-loop system in the \(z\) domain as \(\mathbf{Z}(z)\bm{q}(z)=\bm{f}_{\mathrm{red}}(z)\), with
\begin{equation}
	\begin{aligned}
		\mathbf{Z}(z)={}&
		\frac{(z-1)^{2}}{z\Delta t^{2}}
		\left(\mathbf{M}_{\mathrm{red}}
		+\tfrac{1}{2}\Delta t\,\mathbf{C}_{\mathrm{red}}
		+\tfrac{1}{4}\Delta t^{2}\mathbf{K}_{\mathrm{red}}\right)
		+\frac{z-1}{z\Delta t}\mathbf{C}_{0}+\mathbf{K}_{0}\\
		&+\sum_{\ell=1}^{n_{\mathrm{a}}}z^{-n_{\ell}}
		\left[
		\frac{z-1}{z\Delta t}\mathbf{C}_{\ell}+\mathbf{K}_{\ell}
		\right]
		+\mathbf{S}_{\mathrm{a}}^{\mathsf{T}}\mathbf{H}_{\mathrm{a}}(z)
		\left[
		\mathbf{K}_{x}+\frac{z-1}{z\Delta t}\mathbf{K}_{v}
		\right]\mathbf{S}_{\mathrm{a}},
	\end{aligned}
	\label{eq:z-matrix}
\end{equation}
\rev{where} \(\mathbf{H}_{\mathrm{a}}(z)\) is the diagonal matrix whose \(\ell\)th entry is the delay factor \(z^{-n_{\ell}}\) of the \(\ell\)th channel. The poles of the closed-loop system are the roots of
\begin{equation}
	\det\mathbf{Z}(z)=0.
	\label{eq:characteristic-equation}
\end{equation}
The negative powers of \(z\) in Eq. \eqref{eq:z-matrix} are cleared by multiplying Eq. \eqref{eq:characteristic-equation} by a sufficient power of \(z\), and the roots introduced at the origin by this operation are discarded.

The continuous and \rev{the} discrete planes are related by \(z=\mathrm{e}^{s\Delta t}\), \rev{and} the left half of the \(s\) plane \rev{therefore maps} onto the interior of the unit circle of the \(z\) plane. The largest modulus among the roots of Eq. \eqref{eq:characteristic-equation} is written as \(\rho_{\mathrm{cl}}\). The system is stable when \(\rho_{\mathrm{cl}}<1\), unstable when \(\rho_{\mathrm{cl}}>1\), and on the stability boundary when \(\rho_{\mathrm{cl}}=1\).

For a given condensed model, \(\rho_{\mathrm{cl}}\) depends on the delay of each actuator through the factors \(z^{-n_{\ell}}\). The stability domain is obtained by evaluating \(\rho_{\mathrm{cl}}\) over combinations of the integer delays \(n_{\ell}\) and by locating the contour on which \(\rho_{\mathrm{cl}}\) equals unity. The delays are expressed in multiples of the sampling period, and the resulting domain gives the combinations of actuator delays for which the RTHS system remains stable.




\section{Benchmark structure and analysis setup} \label{sec4}

The formulations of Sections \ref{sec2} and \ref{sec3} are \rev{applied} to a benchmark frame \rev{whose} physical substructure contains more interface coordinates than the transfer system can impose independently. The study is a virtual RTHS, in which both substructures are simulated and the two-channel actuation limit is imposed as a modelling constraint. \rev{The accuracy and the stability of a condensed model depend on the extent of the physical substructure and on the choice of the actuated coordinates, and two substructure divisions with different condensation ratios are therefore examined.} In this section, the benchmark structure and its finite element idealization are described first. The two substructure divisions are then defined together with the retained and condensed coordinates of each substructure. The excitations, the controller and delay settings, and the indices \rev{adopted for the modelling accuracy and the stability} are given last.

\subsection{Benchmark structure and finite element idealization}
\label{sec:benchmark}

The reference structure is a simplified planar model built from the geometry and the material data of the multi-axial RTHS benchmark frame of Condori Uribe et al. \cite{Condori2023}\,\rev{[DOI \nolinkurl{10.3389/fbuil.2023.1270996}]}. The three bays are 762 mm wide and the three storeys are 635 mm high. The columns are W5\(\times\)16 sections of A992 Gr.50 steel and the beams are a custom I section of A36 steel. \rev{The geometry and the section dimensions are shown in Fig. \ref{fig:benchmark-geometry}, and the material properties are listed in Table \ref{tab:materials}.}

\begin{figure}[htbp]
	\centering
	% Selected option: Figure 3-C. Compare with TU thesis Fig. 2-3 (PDF p.28).
	\includegraphics[width=0.89\textwidth]{figure/selected_v2/fig03_benchmark_geometry_optionC.png}
	\caption{Geometry and member sections of the benchmark frame.}
	\label{fig:benchmark-geometry}
\end{figure}

\begin{table}[htbp]
	\centering
	\caption{Material properties of the benchmark frame.}
	\label{tab:materials}
	\begin{tabular}{lccccc}
		\toprule
		Steel grade & \(f_{\mathrm{y}}\) (MPa) & \(f_{\mathrm{u}}\) (MPa) & \(E\) (GPa) & \(\nu\) & \(\rho\) (kg/m\(^{3}\))\\
		\midrule
		A36 (beams)      & 250 & 400--550 & 200 & 0.3 & 7850\\
		A992 Gr.50 (columns) & 345 & 450--550 & 200 & 0.3 & 7850\\
		\bottomrule
	\end{tabular}
\end{table}

In the finite element idealization, each node of the frame carries two in-plane translations and one rotation about the out-of-plane axis. The frame responds mainly in the horizontal direction under a base excitation\rev{. The axial deformation of the members is therefore neglected}, the vertical displacement is taken as zero, and the horizontal displacement is assumed \rev{identical for} all nodes of a storey. The idealized model of Fig. \ref{fig:dof-idealization} \rev{has} fifteen DOFs, which are one horizontal displacement per storey and one rotation at each of the four nodes of each storey. \rev{Among them}, \(\psi_{1}\), \(\psi_{6}\) and \(\psi_{11}\) \rev{are} the horizontal displacements of the first, second and third storey, and the remaining coordinates are the nodal rotations of the corresponding storey.

\begin{figure}[htbp]
	\centering
	% Selected option: Figure 4-D.
	% Panel (a): TU thesis Fig. 2-4 (PDF p.29); panel (b): Fig. 2-5 (PDF p.30).
	\includegraphics[width=0.77\textwidth]{figure/selected_v2/fig04b_dof_idealized_optionD.png}
	\caption{Degrees of freedom of the idealized model.}
	\label{fig:dof-idealization}
\end{figure}

The full-order equation of motion is Eq. \eqref{eq:full-order-eom}. The damping matrix is proportional to the mass and stiffness matrices, \(\mathbf{C}=a_{0}\mathbf{M}+a_{1}\mathbf{K}\), and the coefficients \(a_{0}\) and \(a_{1}\) are obtained from a damping ratio of 5\% at the first two modes. The first five natural frequencies of the idealized model are 2.7 Hz, 9.1 Hz, 18.0 Hz, 22.1 Hz and 23.6 Hz. The first five modes account for more than 90\% of the effective modal mass and therefore characterize the dominant seismic response of the structure.

\subsection{Substructure division and selection of retained DOFs}
\label{sec:division}

The position of the boundary between the physical and the numerical substructures affects both the dynamic response and the accuracy of the condensation. Two divisions are compared and both are shown in Fig. \ref{fig:division}. In each of them the outer bay of the frame is implemented as the physical substructure and the remaining part forms the numerical substructure. The outer bay is selected because a division at an inner bay would split the numerical substructure into two separate parts.

In division I the physical substructure covers the first two storeys of the outer bay together with the connected beams and columns. The lower storeys carry the largest storey shear of the first two modes, and the interface is close to the actuator locations, which shortens the loading path. The physical substructure holds about 40\% of the DOFs of the whole structure.

In division II the physical substructure covers all three storeys of the outer bay. More DOFs are kept at the interface, \rev{which represents the modal coupling between the storeys more completely, and the physical substructure holds} about 60\% of the DOFs of the whole structure.

\rev{Three requirements govern the selection of the retained coordinates. The coordinates carrying the dominant response remain explicit, the condensed response is recoverable through the relations of Section \ref{sec2}, and the reduction is large enough to be of practical use.} The horizontal storey displacements govern the global response of a frame under a base excitation and carry most of the energy of the low-order modes. The nodal rotations mainly represent local bending and are coupled to the horizontal displacements through the off-diagonal terms of the stiffness matrix, \rev{and they are condensed accordingly. Condensing all rotations would shift the higher modes, and several rotations of the numerical substructure are therefore kept.}

Both divisions are limited to two independent actuation channels. Division I directly actuates both horizontal coordinates of its physical substructure\rev{. Division II has} three horizontal coordinates under the same two-channel constraint, \rev{and} the middle-storey coordinate \(\psi_{6}\) is \rev{therefore} condensed and reconstructed from the retained coordinates. Division II is \rev{the more demanding case, since a coordinate carrying} a substantial share of the storey response is no longer imposed directly. \rev{The retained horizontal coordinates of the physical substructure are the coordinates imposed by the two actuators in both divisions. The retained set and the actuated set thus coincide, and the same assignment is adopted in the accuracy analysis and in the stability analysis.}

Three rotations of the numerical substructure are retained as internal master coordinates, which limits the shift of its higher modes. \rev{They leave the interface coordinates condensed by the physical substructure outside the coupling set, where the compatibility is not restored.} The retained and condensed coordinates of both substructures are listed in Table \ref{tab:dof-sets}\rev{ with} the coordinate indices of Fig. \ref{fig:dof-idealization}, and a coordinate \rev{lying} on the interface appears in both substructures.

\begin{figure}[htbp]
	\centering
	% Selected option: Figure 5-D.
	% Panel (a): TU thesis Fig. 3-1; panel (b): Fig. 3-2; both on PDF p.37.
	\includegraphics[width=0.89\textwidth]{figure/selected_v2/fig05a_divisionI_concept_optionD.png}
	\par\vspace{1.5mm}
	\includegraphics[width=0.89\textwidth]{figure/selected_v2/fig05b_divisionII_concept_optionD.png}
	\caption{Substructure divisions and selection of the retained coordinates. (a) Division I. (b) Division II.}
	\label{fig:division}
\end{figure}

\begin{table}[htbp]
	\centering
	\caption{Retained and condensed coordinates of the two substructure divisions.}
	\label{tab:dof-sets}
	\begin{tabular}{llll}
		\toprule
		Division & Substructure & Retained coordinates & Condensed coordinates\\
		\midrule
		I  & Numerical & \(\psi_{1},\psi_{4},\psi_{6},\psi_{9},\psi_{11},\psi_{14}\)
		& \(\psi_{3},\psi_{5},\psi_{7},\psi_{8},\psi_{10},\psi_{12},\psi_{13},\psi_{15}\)\\
		I  & Physical  & \(\psi_{1},\psi_{6}\)
		& \(\psi_{2},\psi_{3},\psi_{7},\psi_{8}\)\\
		\midrule
		II & Numerical & \(\psi_{1},\psi_{4},\psi_{6},\psi_{9},\psi_{11},\psi_{14}\)
		& \(\psi_{3},\psi_{5},\psi_{8},\psi_{10},\psi_{13},\psi_{15}\)\\
		II & Physical  & \(\psi_{1},\psi_{11}\)
		& \(\psi_{2},\psi_{3},\psi_{6},\psi_{7},\psi_{8},\psi_{12},\psi_{13}\)\\
		\bottomrule
	\end{tabular}
\end{table}

The three fixed-interface modes of lowest frequency are retained for each substructure in the Craig--Bampton models, which gives the two divisions the same Craig--Bampton order. Assembly through the two interface coordinates of Table \ref{tab:dof-sets} \rev{gives} a Guyan model with 6 generalized coordinates and a Craig--Bampton model with 12 \rev{in both divisions}. The two divisions are therefore compared at the same reduced model size, \rev{and the differences between them do not come from} the order of the reduced model.

\subsection{Simulation setup and evaluation indices}
\label{sec:setup}

The full-order model, the Guyan condensed model and the Craig--Bampton condensed model are implemented in Simulink. \rev{The coupling is idealized in the accuracy study, where the controller dynamics, the measurement noise and the delay are absent and} the displacement computed by the numerical substructure reaches the physical substructure without error. \rev{Any difference between the full-order and the condensed response then comes from the reduction and from the reduced-interface coupling.} The ground acceleration is multiplied by the mass matrix to give the equivalent seismic force. The displacement of the numerical substructure is passed to the physical substructure, which returns the restoring force at the interface. The full-order model is coupled at the complete set of interface coordinates, whereas the two condensed models are coupled only at the interface coordinates \rev{retained by both}.

Two excitations are \rev{adopted}. The first is the El Centro record scaled by a factor of 0.40, which keeps the interstorey drift \rev{moderate} and the member forces below yield, \rev{and the linear model of Eq. \eqref{eq:full-order-eom} therefore represents the frame}. The second is a chirp signal whose frequency increases linearly from 0.1 Hz to 10 Hz over 40 s. \rev{The sweep covers the first two natural frequencies of the benchmark and follows the accuracy of the condensed model as the excitation frequency approaches and then exceeds the fundamental frequency.}

The stability analysis \rev{adopts} the closed loop of Section \ref{sec3}. The sampling period is \(\Delta t=1/1024\) s and \rev{equals} the integration time step. Two actuators drive the retained horizontal coordinates of the physical substructure. \rev{Actuator 1 drives \(\psi_{1}\) with a delay \(\tau_{1}\) in both divisions, while actuator 2 drives \(\psi_{6}\) in division I and \(\psi_{11}\) in division II, both with a delay \(\tau_{2}\).} The two delays are treated as independent, \rev{which represents channels of different equivalent pure delay}. The weighting matrices of the regulator are \(\mathbf{Q}=\mathrm{diag}(10^{6},10^{6},10^{4},10^{4})\) and \(\mathbf{R}=\mathrm{diag}(10^{-2},10^{-2})\), \rev{where} the first two entries of \(\mathbf{Q}\) weight the displacements of the two actuated coordinates and the last two weight their velocities.

Three indices are \rev{adopted for} the accuracy of the condensed models. The first is the relative error of the natural frequencies, \(e_{f,i}=|f_{i}^{\mathrm{full}}-f_{i}^{\mathrm{red}}|/f_{i}^{\mathrm{full}}\times100\%\), \rev{where} \(f_{i}^{\mathrm{full}}\) and \(f_{i}^{\mathrm{red}}\) are the \(i\)th natural frequency of the full-order and of the condensed model. The second is the modal assurance criterion,
\begin{equation}
	\mathrm{MAC}_{i}=
	\frac{\left[\left(\bm{\varphi}_{i}^{\mathrm{full}}\right)^{\mathsf{T}}
		\bm{\varphi}_{i}^{\mathrm{red}}\right]^{2}}
	{\left\|\bm{\varphi}_{i}^{\mathrm{full}}\right\|^{2}
		\left\|\bm{\varphi}_{i}^{\mathrm{red}}\right\|^{2}},
	\label{eq:mac}
\end{equation}
where \(\bm{\varphi}_{i}^{\mathrm{red}}\) is the \(i\)th mode shape of the condensed model and \(\bm{\varphi}_{i}^{\mathrm{full}}\) is the \(i\)th mode shape of the full-order model restricted to the retained coordinates, \rev{which places the two vectors on the same coordinate set. The retained-coordinate part of the reduced mode shape is taken for the Craig--Bampton models. A value close to unity indicates preserved mode shape components on the retained coordinates.} The third is the normalized root mean square error of the time histories,
\begin{equation}
	\mathrm{NRMSE}=
	\frac{\displaystyle\sqrt{\frac{1}{T_{\mathrm{d}}}\int_{0}^{T_{\mathrm{d}}}
			\left[x^{\mathrm{full}}(t)-x^{\mathrm{red}}(t)\right]^{2}\mathrm{d}t}}
	{\max\left(x^{\mathrm{full}}\right)-\min\left(x^{\mathrm{full}}\right)}
	\times100\%,
	\label{eq:nrmse}
\end{equation}
where \(x^{\mathrm{full}}\) and \(x^{\mathrm{red}}\) are the horizontal storey displacements of the full-order and of the condensed model\rev{, and} \(T_{\mathrm{d}}\) is the duration of the response. \rev{For a frequency band, the numerator is evaluated over the corresponding time interval and divided by the duration of that interval, which places bands of different length on the same root mean square basis. The denominator is the peak-to-peak value of the full-order response over the whole record and is common to every band, and a band of small response amplitude therefore returns a small NRMSE.}

\rev{For the chirp excitation, the} NRMSE is evaluated separately over five frequency bands, which are 0.1 Hz to \(0.7f_{1}\), \(0.7f_{1}\) to \(1.3f_{1}\), \(1.3f_{1}\) to \(2f_{1}\), \(2f_{1}\) to \(3f_{1}\), and \(3f_{1}\) to 10 Hz, \rev{where} \(f_{1}\) is the fundamental frequency of the structure. With \(f_{1}=2.7\) Hz these bands correspond to 0.1 to 1.9 Hz, 1.9 to 3.5 Hz, 3.5 to 5.4 Hz, 5.4 to 8.1 Hz and 8.1 to 10 Hz. The instantaneous frequency of the sweep is \(f(t)=0.1+0.2475t\)\rev{. The band containing the fundamental frequency therefore spans the interval from about 7.3 s to 13.7 s, and} the sweep crosses \(f_{1}\) at about 10.5 s. The bands are centred on the fundamental frequency, where the condensation error is most sensitive to the excitation frequency. The response within a band \rev{carries} the decay of the preceding band, \rev{and the band values accordingly trace the growth of the error with the excitation frequency}.

\rev{The effect of condensation on the stability is quantified through a modal redistribution ratio, which measures the low-order modal content moved onto the retained coordinates. With the mode shapes of the full-order model normalized with respect to the mass matrix, the modal share of coordinate \(j\) is expressed as}
\begin{equation}
	\begin{gathered}
		\gamma_{i}=
		\frac{\bm{\varphi}_{i}^{\mathsf{T}}\mathbf{M}\bm{\Gamma}}
		{\bm{\varphi}_{i}^{\mathsf{T}}\mathbf{M}\bm{\varphi}_{i}},
		\qquad
		p_{i}=\frac{\gamma_{i}^{2}}{\displaystyle\sum_{k=1}^{m}\gamma_{k}^{2}},\\[2pt]
		E_{j}=\sum_{i=1}^{m}p_{i}\,
		\frac{m_{j}\varphi_{j,i}^{2}}{\displaystyle\sum_{l=1}^{n}m_{l}\varphi_{l,i}^{2}},
	\end{gathered}
	\label{eq:dof-energy}
\end{equation}
where \rev{\(\gamma_{i}\) is the modal participation factor of the \(i\)th mode, \(p_{i}\) is the corresponding modal weight, \(E_{j}\) is the modal share of coordinate \(j\),} \(\bm{\varphi}_{i}\) is the \(i\)th mode shape and \(\varphi_{j,i}\) is its \(j\)th entry, \(\bm{\Gamma}\) is the influence vector of Eq. \eqref{eq:full-order-eom}, \(m_{j}\) is the diagonal entry of the mass matrix associated with coordinate \(j\), \(n\) is the number of DOFs of the full-order model, and \(m\) is the number of modes considered, taken as \(m=5\) for the frame of Section \ref{sec:benchmark}. The mode shapes of a condensed model are expanded to the full coordinate set through its own transformation matrix before Eq. \eqref{eq:dof-energy} is applied, \rev{which evaluates the two shares on the same coordinates. The modal redistribution ratio is then obtained as}
\begin{equation}
	\Delta E=\sum_{k=1}^{n_{\mathrm{a}}}
	\frac{E^{\mathrm{red}}(d_{k})-E^{\mathrm{full}}(d_{k})}{E^{\mathrm{full}}(d_{k})}.
	\label{eq:energy-change}
\end{equation}
\rev{where \(d_{k}\) is the \(k\)th actuated coordinate of the physical substructure, and \(E^{\mathrm{full}}\) and \(E^{\mathrm{red}}\) are the modal shares before and after condensation. Equation \eqref{eq:energy-change} sums the relative change of the modal share of each actuated coordinate. A larger \(\Delta E\) therefore corresponds to a larger part of the modal content of the condensed coordinates carried by the coordinates that the actuators drive, and \(\Delta E\) measures the response exposed to the actuator delay through the propagation of Section \ref{sec:multiple-delay}. The index is built from the diagonal entries of the mass matrix and is applied to comparisons on the same physical coordinate set.}




\section{Results and discussion} \label{sec5}

The two condensed models of each division are \rev{compared} with the full-order reference model of Section \ref{sec:benchmark}\rev{, following} the order in which condensation acts on an RTHS. The modal properties of the reduced models are examined first, \rev{since} a shift of the natural frequencies changes the response at every excitation level. The response accuracy under the two excitations is evaluated next\rev{ with} the idealized coupling\rev{, where} condensation is the only source of error. The stability domains of the closed loop are then computed for pairs of actuator delays. \rev{The modal redistribution ratio finally relates the two sets of results to the modal content moved onto the actuated coordinates by condensation.}

\subsection{Modal accuracy}
\label{sec:modal-accuracy}

The relative errors of the first two natural frequencies and the corresponding MAC values are listed in Table \ref{tab:modal}. The chirp excitation reaches 10 Hz, \rev{and the first two modes at 2.7 Hz and 9.1 Hz are therefore the modes within the excitation band, to which the comparison is restricted}. Craig--Bampton condensation reproduces the natural frequencies of the full-order model in both divisions, with a largest error of 0.834\%. Guyan condensation is accurate for the first mode of division I, where the error is 0.648\%, \rev{while} the error of the second mode \rev{reaches} 5.411\%. \rev{Both modes are affected in division II}, with errors of 13.040\% and 24.575\%.

\begin{table}[htbp]
	\centering
	\caption{Relative error of the natural frequencies and MAC values of the condensed models.}
	\label{tab:modal}
	\begin{tabular}{llcccc}
		\toprule
		& & \multicolumn{2}{c}{Frequency error (\%)} & \multicolumn{2}{c}{MAC}\\
		\cmidrule(lr){3-4}\cmidrule(lr){5-6}
		Division & Mode & Guyan & Craig--Bampton & Guyan & Craig--Bampton\\
		\midrule
		I  & 1 & 0.648  & 0.003 & 1.0000 & 1.0000\\
		I  & 2 & 5.411  & 0.284 & 0.9998 & 1.0000\\
		\midrule
		II & 1 & 13.040 & 0.007 & 0.9962 & 1.0000\\
		II & 2 & 24.575 & 0.834 & 0.9255 & 0.9998\\
		\bottomrule
	\end{tabular}
\end{table}

The two divisions differ in that division II condenses a horizontal storey coordinate inside a larger physical substructure. Once a horizontal storey coordinate \rev{leaves} the retained set, the reduced stiffness of the physical substructure is built from a static relation between the bottom and the top storey alone, and the inertia associated with the middle storey is redistributed over the two retained coordinates through Eq. \eqref{eq:guyan-transformation}. The natural frequencies of the assembled model shift accordingly. Both condensations restrict the model to a subspace of the full-order one, \rev{and the condensed frequencies therefore lie above the full-order frequencies. Craig--Bampton condensation escapes the same shift, since} the fixed-interface modes retain part of the internal inertia of the physical substructure.

The MAC values behave differently. All four models keep the MAC above 0.92, and the lowest value\rev{ of 0.9255 belongs} to the Guyan model of division II\rev{, whose} second natural frequency is \rev{in} error by 24.575\%. The mode shape components on the retained coordinates are therefore preserved where the frequencies are not. \rev{The Guyan transformation builds} the condensed coordinates from a static deformation pattern that resembles the low-order mode shapes \rev{and carries} no information \rev{on} the inertia \rev{setting} the frequencies. A MAC evaluated on the retained coordinates \rev{alone is accordingly insufficient} to accept a condensed model, and the frequency error is the modal quantity \rev{to be monitored} alongside it.

\subsection{Response accuracy under earthquake and chirp excitation}
\label{sec:response-accuracy}

Figs. \ref{fig:eq-div1} and \ref{fig:eq-div2} compare the horizontal storey displacements of the two condensed models with those of the full-order model under the El Centro record. The first and the third storey are shown. The first storey is actuated in both divisions and allows the same coordinate to be \rev{followed} throughout, and the third storey carries the largest lateral response of the low-order modes. \rev{The response of the second storey lies between the two and is omitted.} In division I both condensed models follow the reference over the whole record, and the deviation of the Guyan model \rev{appears} only in the enlarged window between 10 s and 11 s. In division II the Guyan response departs from the reference over the same window in both amplitude and shape, whereas the Craig--Bampton response stays close to it. The window between 21.5 s and 22.5 s\rev{ carries a weaker excitation and is reproduced by} both models in both divisions, \rev{and} the error of the Guyan model is \rev{therefore} not uniform in time.

\begin{figure}[htbp]
	\centering
	% Calculation-result plot: source/results mapping in figure/results_v2.
	% Compare with TU thesis Fig. 3-6 (PDF p.42) and Fig. 3-8 (PDF p.43).
	\includegraphics[width=0.98\textwidth]{figure/results_v2/PDF/fig06_eq_div1.pdf}
	\caption{Horizontal storey displacement of division I under the El Centro record. (a) First storey. (b) Third storey.}
	\label{fig:eq-div1}
\end{figure}

\begin{figure}[htbp]
	\centering
	% Calculation-result plot: source/results mapping in figure/results_v2.
	% Compare with TU thesis Fig. 3-9 (PDF p.44) and Fig. 3-10 (PDF p.45).
	\includegraphics[width=0.98\textwidth]{figure/results_v2/PDF/fig07_eq_div2.pdf}
	\caption{Horizontal storey displacement of division II under the El Centro record. (a) First storey. (b) Third storey.}
	\label{fig:eq-div2}
\end{figure}

Figs. \ref{fig:chirp-div1} and \ref{fig:chirp-div2} show the response under the chirp excitation, in which the instantaneous frequency rises linearly from 0.1 Hz to 10 Hz over the 40 s of the record. The sweep crosses the fundamental frequency at about 10.5 s and reaches three times the fundamental frequency at about 32 s. In the low-frequency part of the sweep both condensed models reproduce the reference in both divisions.

In division I the Guyan response deviates from the reference around the first resonance peak\rev{ near} 13 s, \rev{and} the deviation remains small. The resonant response builds up over several cycles, \rev{and this peak therefore lags} the crossing of the fundamental frequency at about 10.5 s. \rev{The deviation grows} again towards the end of the sweep\rev{ near} 38 s, where the instantaneous frequency is about 9.5 Hz and the excitation \rev{passes} the second mode of the structure. The second natural frequency of this model is in error by 5.411\%, \rev{which places the second resonance at the wrong frequency and turns the deviation into an amplitude and phase error over that window}. In division II the Guyan response is \rev{inaccurate over the whole} resonance band, and its peaks are shifted in time against those of the reference, which is the time-domain signature of the frequency shift \rev{of} Section \ref{sec:modal-accuracy}. Near the end of the sweep\rev{,} the Guyan response of division II falls to less than half of the reference at the first storey and rises to nearly twice the reference at the third storey. The Craig--Bampton response follows the reference over the whole sweep in both divisions.

\begin{figure}[htbp]
	\centering
	% Calculation-result plot: source/results mapping in figure/results_v2.
	% Compare with unique thesis Fig. 3-11 (PDF p.46) and Fig. 3-13 (p.47--48).
	\includegraphics[width=0.98\textwidth]{figure/results_v2/PDF/fig08_chirp_div1.pdf}
	\caption{Horizontal storey displacement of division I under the chirp excitation. (a) First storey. (b) Third storey.}
	\label{fig:chirp-div1}
\end{figure}

\begin{figure}[htbp]
	\centering
	% Calculation-result plot: source/results mapping in figure/results_v2.
	% Compare with unique thesis Fig. 3-14 (PDF p.48 lower) and Fig. 3-15 (PDF p.49).
	\includegraphics[width=0.98\textwidth]{figure/results_v2/PDF/fig09_chirp_div2.pdf}
	\caption{Horizontal storey displacement of division II under the chirp excitation. (a) First storey. (b) Third storey.}
	\label{fig:chirp-div2}
\end{figure}

Table \ref{tab:nrmse} gives the NRMSE of the first-storey horizontal displacement. Under the El Centro record the Guyan error is 0.758\% in division I and 10.936\% in division II, whereas the Craig--Bampton error stays below 0.03\% in both. The band decomposition of the chirp response \rev{locates the origin of} the Guyan error. Below the fundamental frequency the Guyan error is 0.030\% in division I and 2.996\% in division II. \rev{It rises} to 1.995\% and 19.059\% \rev{in the band containing the fundamental frequency, the latter being the largest value of the table}. Above the fundamental frequency the division I error falls back to 0.076\% in the band from 5.4 Hz to 8.1 Hz and \rev{rises} again to 0.745\% in the highest band\rev{, whose small absolute response places the corresponding NRMSE below the relative deviation of Fig. \ref{fig:chirp-div1} near 38 s}. In division II the error stays above 2.8\% in every band. The Craig--Bampton error remains below 0.06\% in every band of both divisions.

\begin{table}[htbp]
	\centering
	\caption{NRMSE of the first-storey horizontal displacement (\%).}
	\label{tab:nrmse}
	\begin{tabular}{lcccc}
		\toprule
		& \multicolumn{2}{c}{Division I} & \multicolumn{2}{c}{Division II}\\
		\cmidrule(lr){2-3}\cmidrule(lr){4-5}
		Excitation & Guyan & Craig--Bampton & Guyan & Craig--Bampton\\
		\midrule
		El Centro           & 0.758 & 0.006    & 10.936 & 0.027\\
		\midrule
		Chirp, 0.1--1.9 Hz  & 0.030 & $<$0.001 & 2.996  & 0.004\\
		Chirp, 1.9--3.5 Hz  & 1.995 & 0.008    & 19.059 & 0.056\\
		Chirp, 3.5--5.4 Hz  & 0.623 & 0.003    & 14.876 & 0.030\\
		Chirp, 5.4--8.1 Hz  & 0.076 & 0.003    & 7.498  & 0.014\\
		Chirp, 8.1--10 Hz   & 0.745 & 0.051    & 2.813  & 0.007\\
		\bottomrule
	\end{tabular}
\end{table}

Two observations follow. The accuracy of a condensed model depends on the excitation frequency, and the band around the fundamental frequency is where the two methods differ most. The two divisions \rev{produce} assembled models of the same size \rev{and differ} by more than an order of magnitude in the Guyan error, \rev{and the order of the reduced model therefore does not explain the accuracy by itself}. Division II \rev{enlarges} the physical substructure and condenses the middle-storey horizontal coordinate\rev{ at the same time, and the combination of the two makes} it the more demanding case.

\subsection{Stability domains under multiple actuator delays}
\label{sec:stability-domains}

The stability domains obtained from the pole criterion of Eq. \eqref{eq:characteristic-equation} are shown in Fig. \ref{fig:stability-domain}. The two axes are the integer delays of the two channels\rev{ expressed} in multiples of the sampling period. \rev{One step corresponds to about 0.98 ms with \(\Delta t=1/1024\) s, and the axes therefore read approximately as milliseconds over the range plotted}. The delays are scanned over integer pairs, and each boundary is the largest stable \(n_{2}\) \rev{obtained} for each \(n_{1}\). Every pair between a boundary and the origin \rev{is} stable in the scan. The Original boundary \rev{belongs to the unreduced fifteen-DOF model, with} the two delays applied to the same two coordinates and \rev{the} remaining interface coordinates coupled without delay. It is \rev{an} upper bound obtained with an ideal interface\rev{ rather than} a model \rev{under} the same two-channel constraint as the condensed ones.

\begin{figure}[htbp]
	\centering
	% Historical stability-boundary snapshot, plotting-level reproduction only.
	% Compare with TU thesis Fig. 4-4 and Fig. 4-5 (both on PDF p.70).
	\includegraphics[width=0.98\textwidth]{figure/results_v2/PDF/fig10_stability_domain.pdf}
	\caption{Stability domains in the plane of the two actuator delays. (a) Division I. (b) Division II.}
	\label{fig:stability-domain}
\end{figure}

Both condensations reduce the stability domain relative to the unreduced model, and in both divisions the ordering of the three boundaries is the same. The unreduced model has the largest domain, the Craig--Bampton model comes next, and the Guyan model has the smallest. The contraction is larger in division II than in division I for both methods.

The recovery relations of Section \ref{sec:multiple-delay} account for the ordering. \rev{The reconstructed response of the condensed coordinates is built from delayed quantities alone under Guyan condensation, which exposes the whole internal response of each substructure to the actuator delay.} Under Craig--Bampton condensation the fixed-interface modal coordinates carry no explicit delay operator and reach the delayed channels only through their coupling to the retained coordinates, \rev{and} the delay \rev{therefore enters} a smaller part of the model. The two models differ in basis and in order, at twelve generalized coordinates against six, and \rev{the separation of the two effects rests on the number of retained fixed-interface modes}.

The stability domain of the unreduced model is itself smaller in division II than in division I. The two divisions differ \rev{in} the extent of the physical substructure and in the position of the second actuator, \rev{and} the combined change lowers the delay margin before any condensation is applied\rev{. Condensation adds to this loss, and} the spacing between the boundaries within one panel is produced by the condensation together with the regulator \rev{supplied by} each condensed model.

\subsection{\rev{Modal redistribution ratio} and applicability of the two methods}
\label{sec:energy-applicability}

The \rev{modal redistribution ratio} of Eq. \eqref{eq:energy-change} is listed in Table \ref{tab:energy}. Within each division the three boundaries of Fig. \ref{fig:stability-domain} are nested, \rev{and} the two condensed models of that division \rev{are therefore ordered} by set inclusion. \rev{The Craig--Bampton value is the smaller of the two in both divisions} and the Craig--Bampton domain is the larger. \rev{One method compared across} the two divisions gives the same correspondence, since both the \rev{modal redistribution ratio} and the contraction of the domain are larger in division II.

\begin{table}[htbp]
	\centering
	\caption{\rev{Modal redistribution ratio} of the four condensed models.}
	\label{tab:energy}
	\begin{tabular}{lcc}
		\toprule
		Division & Guyan & Craig--Bampton\\
		\midrule
		I  & 0.3065 & 0.1879\\
		II & 0.3927 & 0.2021\\
		\bottomrule
	\end{tabular}
\end{table}

\rev{Table \ref{tab:energy} is the quantitative counterpart of the mechanism of Section \ref{sec:multiple-delay}. A large modal redistribution ratio corresponds to a large share of the modal energy originally distributed over the condensed coordinates moved onto the retained coordinates. The retained coordinates of the physical substructure are the coordinates driven by the actuators, and the same quantity therefore measures the response of the reduced model carried by delayed channels. The ratio is obtained from the modal solutions of the full-order and of the condensed model, and it is available as soon as the condensed model is built, ahead of any closed-loop delay scan.}

\rev{The two models with a modal redistribution ratio above 0.3 are the two models whose stability domain contracts most, and both are Guyan models. The ratio ranks candidate condensed models ahead of the delay scan, and the four models of Table \ref{tab:energy} support a ranking rather than a general threshold.}

\rev{The accuracy and the stability results together indicate the range of application of each method.} Craig--Bampton condensation keeps the frequency error below 1\% and the NRMSE below 0.06\% in every case examined, and its stability boundary is the closer of the two to that of the unreduced model\rev{. It is therefore the method for a condensed principal coordinate and for an excitation carrying energy near or above the fundamental frequency. Guyan condensation produces the smallest reduced model, at six generalized coordinates against twelve for Craig--Bampton. In division I its second natural frequency is in error by 5.411\% and its NRMSE reaches 1.995\% in the band around the fundamental frequency, which sets the tolerance required of the application. In division II its frequency error reaches 24.575\% and its NRMSE reaches 19.059\%, more than an order of magnitude above the Craig--Bampton values, and it also gives the smallest stability domain.}

\rev{The three indices} do not carry equal weight in this judgement. The frequency error and the NRMSE of the Guyan model of division II are the largest values \rev{of} Tables \ref{tab:modal} and \ref{tab:nrmse}, whereas the MAC of the same model remains above 0.92. A modal correlation index evaluated on the retained coordinates \rev{alone is therefore insufficient} to accept a condensed model for RTHS, and it \rev{is applied} together with a frequency error and a time-domain error.




\section{Conclusions}\label{sec6}

This study establishes that the reduction basis of an actuation-limited MDOF RTHS determines the path of the actuator delays as well as the approximation of the structural response. Guyan condensation reconstructs every condensed physical DOF from delayed retained DOFs, so its constraint matrix distributes the delayed channels throughout the recovered response. Craig--Bampton condensation adds fixed-interface modal coordinates that are solved inside the reduced model and carry no explicit delay operator. The channel-wise reduced formulation makes this distinction explicit and embeds both transformations in the same discrete-time closed loop. The dominant-pole criterion then provides the stability domain for independent actuator delays. The modal redistribution ratio complements the delay scan by measuring the low-order modal content transferred to the actuator-driven DOFs when the condensed model is formed.

The benchmark results show that dynamic fidelity and delay margin depend on the selected basis and on which DOFs are condensed. Craig--Bampton condensation keeps the errors of the first two natural frequencies below 0.834\% and the first-storey NRMSE below 0.06\% in every frequency band of both divisions. The Guyan error concentrates around the fundamental frequency and reaches 1.995\% in Division I and 19.059\% in Division II, while its second natural-frequency error reaches 24.575\% in Division II. All four reduced models retain a MAC above 0.92, so a MAC evaluated on the retained DOFs alone does not reveal large frequency and time-response errors. Both condensations contract the stability domain relative to the ideal unreduced reference. In both divisions, the Craig--Bampton boundary lies outside the Guyan boundary, and the contraction grows when the middle-storey horizontal DOF is condensed in Division II. This ordering is consistent with the delay-recovery relations.

The modal redistribution ratios give the same model ranking. The values are 0.3065 and 0.3927 for Guyan condensation and 0.1879 and 0.2021 for Craig--Bampton condensation in Divisions I and II. These cases establish the ratio as a comparative screening index rather than a universal acceptance threshold. Guyan condensation gives the smaller model, with six generalized coordinates, and applies when the condensed response follows the retained DOFs through quasi-static deformation and the excitation remains below the fundamental frequency. Craig--Bampton condensation uses twelve generalized coordinates in the present study and is the appropriate choice when a principal response DOF is condensed or when the input contains appreciable energy near or above the fundamental frequency. Model order alone therefore does not determine either accuracy or delay margin. The present linear virtual-RTHS formulation with constant integer channel delays provides the basis for experimental validation with time-varying delays and nonlinear physical substructures.

\begin{thebibliography}{9}
	
	\bibitem{Guyan1965}
	Guyan RJ. Reduction of stiffness and mass matrices. AIAA Journal 1965;3(2):380. \url{https://doi.org/10.2514/3.2874}
	
	\bibitem{CraigBampton1968}
	Craig RR, Bampton MCC. Coupling of substructures for dynamic analyses. AIAA Journal 1968;6(7):1313--1319. \url{https://doi.org/10.2514/3.4741}
	
	\bibitem{Carrion2007}
	Carrion JE, Spencer BF. Model-based strategies for real-time hybrid testing. NSEL Report No.\ NSEL-006, University of Illinois at Urbana-Champaign; 2007.
	
	% Kept for the Introduction, where the fractional-step delay model is reviewed.
	% It is no longer cited in Section 3.
	\bibitem{ChenRicles2008a}
	Chen C, Ricles JM. Stability analysis of SDOF real-time hybrid testing systems with explicit integration algorithms and actuator delay. Earthquake Engineering and Structural Dynamics 2008;37(4):597--613. \url{https://doi.org/10.1002/eqe.775}
	
	\bibitem{Zhu2015}
	Zhu F, Wang JT, Jin F, Chi FD, Gui Y. Stability analysis of MDOF real-time dynamic hybrid testing systems using the discrete-time root locus technique. Earthquake Engineering and Structural Dynamics 2015;44(2):221--241. \url{https://doi.org/10.1002/eqe.2467}
	
	\bibitem{ChenRicles2008b}
	Chen C, Ricles JM. Development of direct integration algorithms for structural dynamics using discrete control theory. Journal of Engineering Mechanics 2008;134(8):676--683. \url{https://doi.org/10.1061/(ASCE)0733-9399(2008)134:8(676)}
	
	\bibitem{Condori2023}
	Condori Uribe JW, Salmeron M, Patino E, Montoya H, Dyke SJ, Silva CE, Maghareh A, Najarian M, Montoya A. Experimental benchmark control problem for multi-axial real-time hybrid simulation. Frontiers in Built Environment 2023;9:1270996. \url{https://doi.org/10.3389/fbuil.2023.1270996}
	
	\bibitem{Krattiger2019}
	Krattiger D, Wu L, Zacharczuk M, Buck M, Kuether RJ, Allen MS, Tiso P, Brake MRW. Interface reduction for Hurty/Craig--Bampton substructured models, review and improvements. Mechanical Systems and Signal Processing 2019;114:579--603. \url{https://doi.org/10.1016/j.ymssp.2018.05.031}
	
	\bibitem{Li2021}
	Li N, Lu XJ, Zhou ZH, et al. Stability of real-time substructure testing considering numerical integration algorithms. Engineering Mechanics 2021;38(11):12--22. [in Chinese]
	
\end{thebibliography}

\end{document}
